\documentclass[journal]{IEEEtran}

\usepackage[T1]{fontenc}
\usepackage[utf8]{inputenc}
\usepackage{amsmath,amssymb,amsfonts}
\usepackage{amsthm}    
\usepackage{algorithmic}
\usepackage{algorithm}
\usepackage{array}
\usepackage{textcomp}
\usepackage{cite}
\usepackage{graphicx}
\usepackage{booktabs}
\usepackage{multirow}
\usepackage{xcolor}
\usepackage{url}
\usepackage{hyperref}
\hypersetup{colorlinks=true, linkcolor=blue, citecolor=blue, urlcolor=blue}

% FIX 2: \newtheorem declarations belong in the preamble, declared ONCE each.
%         They were previously declared three times each inside the document
%         body, which causes LaTeX to throw "already defined" errors.
\newtheorem{definition}{Definition}
\newtheorem{proposition}{Proposition}

\definecolor{ieeeblue}{RGB}{0,114,188}

% FIX 3: \texttt in two-column IEEE layout hyphenates long identifiers.
%         This macro wraps them in \mbox to prevent mid-token line breaks.
\newcommand{\mut}[1]{\mbox{\texttt{#1}}}

% ─────────────────────────────────────────────────────────────────────────────
\begin{document}

\title{Semantic Contract-Based Validation for Detecting Data 
Pipeline Mutations in Tabular\\Machine Learning Datasets}

\author{Muhammad Idrees
and Adnan Amin
\thanks{M. Idrees is with the Institute of Management Sciences (IMSciences), 
Peshawar, Pakistan (e-mail: muhammadidrees9405@gmail.com).}
\thanks{A. Amin is with the Institute of Management Sciences (IMSciences), 
Peshawar, Pakistan (e-mail: adnan.amin@imsciences.edu.pk). 
A. Amin is the corresponding author.}}

\markboth{IEEE Access}%
{M. Idrees and A. Amin: Semantic Contract-Based Validation for Data Pipeline Mutation Detection}

\maketitle

% ── Abstract ─────────────────────────────────────────────────────────────────
\begin{abstract}
Data quality failures in automated machine learning pipelines are among the
most consequential and least visible reliability risks in deployed AI systems:
they corrupt model behaviour without triggering any observable system error.
A critical failure class---the \emph{semantic data mutation}---arises when
incoming data is structurally well-formed and distributionally plausible yet
violates domain-specific rules that govern the meaning of its columns.
Classic examples include a column swap between stock Open and Close prices,
a weekly sales value inflated sevenfold by a mis-applied daily-to-weekly
aggregation, and an all-dates shift that destroys temporal alignment without
altering any marginal distribution.  No existing production tool is designed
to detect this class of failure.

We introduce a \emph{semantic contract framework}: a set of programmatic,
domain-aware detectors that express what the data \emph{means}, not merely
what it \emph{looks like}.  We formalise the notion of a semantic mutation,
define a five-category taxonomy of 28 mutation operators, and implement 25
domain-specific contracts across three publicly available tabular datasets
(Walmart retail sales, AAPL five-year stock prices, UCI household power
consumption).

Evaluation over $28 \times 4 = 112$ mutated trials and 12 clean-baseline
trials demonstrates that the semantic validator achieves a pooled detection
rate of \textbf{96.4\%} (95\% Wilson CI: 91.2\%--98.6\%), compared with
49.1\%, 42.0\%, and 42.9\% for ensemble statistical, KS-drift, and
Great-Expectations baselines respectively (all McNemar $p < 0.0001$).
False-positive rate is \textbf{0.000} across all validators.  Semantic
precision---the fraction of detections attributed to the correct mutation
category---is 96.3\%, with four of five categories at 100\%.  On a held-out partition of nine mutation types
withheld during development, the semantic validator achieves 88.9\%
(95\% CI: 74.7\%--95.6\%), maintaining a $>$36 percentage-point advantage
over all baselines on unseen mutations.  These results establish that
statistical methods are structurally insufficient for semantic mutation
detection, and that domain-aware contracts provide a principled, practical
complement to existing pipeline monitoring tools.
\end{abstract}

\begin{IEEEkeywords}
Data Validation, Machine Learning, MLOps, Data Quality, Mutation Testing,
Semantic Contracts, Tabular Data, Data Pipelines, Distribution Drift,
Pipeline Monitoring
\end{IEEEkeywords}


\section{Introduction}
\label{sec:intro}


\IEEEPARstart{M}{odern} machine learning systems are not static artefacts: they are
continuously retrained on data delivered by automated pipelines, and their
predictive reliability is only as sound as the data flowing through those
pipelines.  Sculley et al.\ \cite{sculley2015} identified data dependencies
as one of the principal sources of hidden technical debt in real-world ML
systems, noting that silent data errors compound over time in ways that
model-level monitoring cannot detect.  Polyzotis et al.\ \cite{polyzotis2018}
subsequently formalised data-lifecycle challenges in production ML at
Google-scale, demonstrating that validation at the data layer---not just at
the model output layer---is a first-class engineering concern.  Despite this
recognition, the tooling ecosystem has not kept pace: existing validators
check schemas and distributions, but neither can express what the data
\emph{means}.

Consider three scenarios that defeat every existing tool.  In the first, a
column-ordering bug in an ETL migration silently swaps the \texttt{Open} and
\texttt{Close} price columns in a financial dataset.  Both columns remain
numeric, both remain within historical range, both pass all schema and
KS-drift checks---yet a momentum model trained on this data learns
inverted intraday signals.  In the second, an aggregation pipeline applies a
daily-to-weekly summation twice, inflating \texttt{Weekly\_Sales} by a
factor of seven.  The inflated values pass range checks because no absolute
threshold was calibrated; they pass distributional checks because the shape
of the distribution is preserved under scaling.  The downstream demand
forecast overstates inventory requirements by $7\times$.  In the third, a
cron job error shifts all timestamps forward by seven days.  Every marginal
distribution is identical to clean data; only the temporal relationship to
external events is destroyed.  A promotion-attribution model trained on this
data attributes holiday spikes to the wrong week in perpetuity.
In each case, the financial or operational consequence of the silent error
dwarfs the cost of the mutation itself---yet the failure is invisible to
every monitoring tool in the practitioner's stack.

These three scenarios share a defining property: they are
\textbf{semantically anomalous but statistically invisible}.  We formalise
this property and call the underlying class of error a
\emph{semantic data mutation}.  Existing tools---including Great
Expectations \cite{greatexpectations2023}, the Kolmogorov-Smirnov drift
detector \cite{kolmogorov1933}, ensemble statistical monitors, and
TensorFlow Data Validation \cite{baylor2017tfx}---operate exclusively on
syntactic or distributional properties of the data.  They cannot, by
construction, detect mutations that preserve those properties.

This paper makes the following four contributions:

% FIX 4: Original used bare "1)" text instead of \item, which breaks the
%         enumerate environment in LaTeX and produces garbled numbered output.
\begin{enumerate}
\item \textbf{Formal taxonomy:} The first systematic five-category taxonomy
  of 28 semantic mutation operators (Scaling, Temporal, Structural, Quality,
  Aggregation), grounded in documented real-world ETL failure patterns across
  three production-grade tabular datasets.

\item \textbf{Semantic contract framework:} Twenty-five domain-aware
  programmatic detectors encoding relational invariants and temporal ordering
  constraints that are, by construction, inexpressible as distributional
  tests.

\item \textbf{Rigorous empirical evaluation:} A four-seed, four-validator
  controlled experiment demonstrating 96.4\% detection rate at 0.000
  false-positive rate---a $>$96\,pp absolute improvement over random and a
  47.3\,pp improvement over the strongest available statistical
  baseline---with McNemar-exact $p < 0.0001$ for all pairwise comparisons.

\item \textbf{Generalisation evidence:} A held-out evaluation on nine
  mutation types entirely withheld during development, demonstrating 88.9\%
  detection on unseen patterns---a $>$36\,pp advantage over all
  baselines---and identifying the sole detection-limit case with principled
  justification.
\end{enumerate}

The remainder of this paper is organised as follows.
Section~\ref{sec:related} surveys existing data validation tools, ML
pipeline testing methods, mutation testing literature, and drift detection
approaches, positioning our contribution relative to each.
Section~\ref{sec:formulation} provides formal definitions of semantic
mutations, detection rate, false positive rate, and semantic precision, and
proves by construction that statistical tests are insufficient.
Section~\ref{sec:framework} presents the mutation taxonomy and contract
design principles.  Section~\ref{sec:setup} details datasets, baselines, and
statistical methods.  Section~\ref{sec:results} presents all experimental
results.  Section~\ref{sec:discussion} discusses the core argument,
limitations, and threats to validity.  Section~\ref{sec:conclusion} concludes
with practical implications and future directions.

% ═════════════════════════════════════════════════════════════════════════════
\section{Related Work}
\label{sec:related}
% ═════════════════════════════════════════════════════════════════════════════

\subsection{Data Validation Tools for ML Pipelines}

The most widely deployed data validation framework is Great Expectations
\cite{greatexpectations2023}, which allows practitioners to define
expectation suites over column statistics (mean, standard deviation,
range, null rate, cardinality) and evaluate them against incoming data.
Amazon Deequ \cite{schelter2018deequ} provides a declarative API for
large-scale data quality verification on Apache Spark, with support for
unit-test-style constraints and automatic constraint suggestion.  Baylor
et al.\ \cite{baylor2017tfx} describe TensorFlow Extended, whose data
validation component compares serving statistics against training statistics
to detect training-serving skew.  The Auto-Validate-by-History system
\cite{schelter2018automating} detects data quality issues in recurring
pipelines by leveraging statistics from historical executions.

The critical limitation shared by all of these tools is that they are
\emph{syntax-first}: they validate that values fall within calibrated
bounds, not that the relationships between values conform to domain
semantics.  A seven-fold sales inflation is invisible to any tool whose
bounds were calibrated on clean data without explicit knowledge that
\texttt{Weekly\_Sales} must be aggregated exactly once per week.  A column
swap is invisible to any tool that checks columns independently.  Polyzotis
et al.\ \cite{polyzotis2018} explicitly identify the gap between statistical
validation and semantic correctness as an open challenge in production ML.
Our framework directly addresses this gap.

\subsection{ML Pipeline Testing}

Breck et al.\ \cite{breck2017} introduced the ML Test Score rubric, which
codifies production-readiness criteria for ML systems including data tests,
feature tests, and model tests.  They identify monitoring of input data
distributions as essential, but acknowledge that semantic correctness checks
are domain-specific and beyond the scope of automated scoring.  Schelter
et al.\ \cite{schelter2018automating} present a system for automating
large-scale data quality verification, demonstrating that constraint
learning from historical data can catch syntactic anomalies at scale.
Shankar et al.\ \cite{shankar2022operationalizing} conducted interview
studies with production ML practitioners and found that data debugging
accounts for the majority of operational effort, with semantic errors
among the hardest to detect.  The continuous validation framework described by Reis and 
Housley~\cite{reis2022} emphasises transformation-rule validation 
as the layer most frequently absent in production pipelines---precisely 
the layer our semantic contracts occupy. M\"{a}kinen et 
al.~\cite{makinen2021mlops} similarly find that data preparation and 
validation constitute the dominant operational burden in production ML 
teams, confirming that the problem this paper addresses is both current 
and widely experienced.

\subsection{Mutation Testing}

Mutation testing originates with DeMillo et al.\ \cite{demillo1978}, who
proposed systematically introducing faults into programs to evaluate test
suite adequacy.  Jia and Harman \cite{jia2010} provide a comprehensive
survey of mutation testing advances, establishing the competent programmer
hypothesis---that real faults are syntactically close to correct
programs---as the theoretical justification for mutation operators.

More recently, mutation testing has been applied to deep learning systems.
Ma et al.\ \cite{ma2018deepmutation} propose mutation operators at the data,
program, and model levels for DL systems.  Humbatova et al.\
\cite{humbatova2021deepcrime} extend this with DeepCrime, defining 35
mutation operators derived from real DL faults, demonstrating that
semantically grounded mutations produce more meaningful test outcomes than
arbitrary perturbations.  Our work applies a parallel methodology to the
data pipeline layer: we define mutations grounded in real ETL failure
patterns (unit confusion, aggregation errors, sensor freeze, column swaps)
rather than in model architecture failures.  This distinction is
fundamental: our mutations attack the \emph{data contract} between pipeline
and model, not the model itself.  No prior mutation testing work targets the
data contract layer of ML pipelines; this paper introduces that paradigm.

\subsection{Distribution Drift Detection}

The Kolmogorov-Smirnov test \cite{kolmogorov1933} remains the most widely
used univariate drift detector in production systems.  Multivariate
extensions include the Population Stability Index \cite{yurdakul2018psi},
Jensen-Shannon divergence \cite{lin1991jsd}, and Wasserstein distance.
Modern monitoring platforms such as Evidently AI and Whylogs provide
ensemble drift scores combining multiple statistical tests.  Our ensemble
baseline implements precisely this approach: KS + JSD + PSI + Pearson
correlation matrix distance + autocorrelation function comparison, with
any-one voting.

The fundamental limitation of all distribution-based methods is established
by construction: mutations that preserve marginal distributions are
undetectable by any test whose power derives solely from distributional
divergence.  A column swap between two highly correlated series may produce
no statistically significant KS statistic on either column.  A timestamp
shift preserves all marginal distributions exactly.  A sensor freeze
producing 120 consecutive identical readings creates a local distributional
anomaly that is diluted to insignificance in a 996,010-row dataset.
Semantic contracts detect all three cases by reasoning about what the
values mean---a capability that is orthogonal to statistical power and
cannot be approximated by adding more statistical tests to an ensemble.

% ═════════════════════════════════════════════════════════════════════════════
\section{Problem Formulation}
\label{sec:formulation}
% ═════════════════════════════════════════════════════════════════════════════

\subsection{Notation and Definitions}

Let $\mathcal{D}$ denote the space of tabular datasets over a fixed schema
$\mathcal{S} = \{c_1, \ldots, c_k\}$ with corresponding domain types
$\{\tau_1, \ldots, \tau_k\}$.  A dataset $D \in \mathcal{D}$ is a finite
sequence of tuples $\langle r_1, \ldots, r_n \rangle$ where each $r_i$
assigns a value in $\tau_j$ to each column $c_j$.

\begin{definition}[Mutation Operator]
A \emph{mutation operator} is a function $\mu : \mathcal{D}
\rightarrow \mathcal{D}$ such that $\mu(D)$ conforms to schema
$\mathcal{S}$ for every $D \in \mathcal{D}$. The output
$D' = \mu(D)$ is called a \emph{mutant} of $D$.
\end{definition}

\begin{definition}[Statistical Indistinguishability]
A mutant $D' = \mu(D)$ is \emph{statistically indistinguishable} from $D$
with respect to a test family $\mathcal{T}$ if
$T(D, D') = \mathrm{pass}$ for all $T \in \mathcal{T}$.  Let
$\mathcal{T}_{\mathrm{stat}}$ denote the family of standard distributional
tests (KS, JSD, PSI, correlation matrix, ACF comparison).
\end{definition}

\begin{definition}[Semantic Contract]
A \emph{semantic contract} is a function
$C : \mathcal{D} \times \mathcal{D} \rightarrow \{0, 1\}$ such that
$C(D_{\mathrm{base}}, D_{\mathrm{test}}) = 0$ indicates a violation of a
domain-specific invariant between $D_{\mathrm{base}}$ and
$D_{\mathrm{test}}$.  A contract returns 1 (satisfied) if and only if the
test dataset conforms to the domain rule encoded by the contract.
\end{definition}

% FIX 5: Added a proper \label to the Semantic Mutation definition so that
%         the subsequent cross-reference uses \ref{def:semantic_mutation}
%         rather than incorrectly referencing the equation label.
\begin{definition}[Semantic Mutation]
\label{def:semantic_mutation}
A mutation operator $\mu$ is \emph{semantic} if there exists a clean dataset
$D \in \mathcal{D}$ such that:
\begin{equation}
T_{\mathrm{stat}}(D,\, \mu(D)) = \mathrm{pass}
\;\;\wedge\;\;
\exists\, C \in \mathcal{C} : C(D,\, \mu(D)) = 0
\label{eq:semantic_mutation}
\end{equation}
where $\mathcal{C}$ is the set of semantic contracts for the domain.
That is, $\mu$ is undetectable by statistical tests but detectable by at
least one semantic contract.
\end{definition}

% FIX 6: Was "Definition~\ref{eq:semantic_mutation}" (wrong — equation label).
%         Corrected to reference the definition label defined above.
Definition~4 formalises the core claim of this paper: a class
of data mutations exists that is, by construction, beyond the
reach of purely statistical validation.  The right-hand conjunct is
satisfiable precisely when the mutation violates a domain invariant that
statistical tests do not encode.

\subsection{Evaluation Metrics}

\begin{definition}[Detection Rate]
Given a validator $V$, a set of $N$ mutated trials $\{(D_i, \mu_i(D_i))\}$,
and the indicator $\mathbf{1}[\cdot]$:
\begin{equation}
\mathrm{DR}(V) = \frac{1}{N} \sum_{i=1}^{N}
\mathbf{1}\!\left[V\!\left(D_i,\, \mu_i(D_i)\right) = \mathrm{fail}\right]
\label{eq:dr}
\end{equation}
\end{definition}

\begin{definition}[False Positive Rate]
Given $M$ clean-baseline trials $\{(D_j, D_j)\}$:
\begin{equation}
\mathrm{FPR}(V) = \frac{1}{M} \sum_{j=1}^{M}
\mathbf{1}\!\left[V\!\left(D_j,\, D_j\right) = \mathrm{fail}\right]
\label{eq:fpr}
\end{equation}
\end{definition}

\begin{definition}[Semantic Precision]
Let $\mathcal{R}$ denote the set of detected mutated trials.  For each
$r \in \mathcal{R}$, let $\kappa(r)$ be the ground-truth category of
$\mu_r$, and let $\hat{\kappa}(r)$ be the category of the highest-confidence
firing detector.  Semantic precision is:
\begin{equation}
\mathrm{SP}(V) = \frac{1}{|\mathcal{R}|}
\sum_{r \in \mathcal{R}} \mathbf{1}\!\left[\hat{\kappa}(r) = \kappa(r)\right]
\label{eq:sp}
\end{equation}
Semantic precision measures not only whether a mutation is detected, but
whether the \emph{correct type of contract} fires---a stronger claim about
the diagnostic utility of the framework.
\end{definition}

\subsection{Why Statistical Tests are Insufficient by Construction}

The KS test and its multivariate generalisations test whether two samples
are drawn from the same marginal distribution.  Their power is zero against
any mutation that is a bijective rearrangement of existing values---including
column permutations, timestamp translations, and any scaling that
redistributes values within the same empirical support.  More precisely:

\begin{proposition}
Let $\mu$ be a column-swap operator that exchanges columns $c_a$ and $c_b$.
If $c_a$ and $c_b$ are drawn from the same marginal distribution,
then $\mathrm{KS}(D_i^{c_a}, \mu(D)_i^{c_a}) = \mathrm{pass}$ with
probability 1 under any sample size.
\end{proposition}

\begin{proof}
A column-swap operator $\mu_{ab}$ satisfies
$\mu_{ab}(D)^{c_a} = D^{c_b}$ and $\mu_{ab}(D)^{c_b} = D^{c_a}$
by definition. If $\mathbb{P}(c_a) = \mathbb{P}(c_b)$, then the
empirical distribution of $\mu_{ab}(D)^{c_a}$ equals that of
$D^{c_a}$, so the two-sample KS statistic
$D_n(D^{c_a}, \mu_{ab}(D)^{c_a}) = 0$ for any $n$, and
the test fails to reject $H_0$ with certainty.
\end{proof}

Similarly, any timestamp shift $\mu_{\delta t}$ satisfies
$\forall c \in \mathcal{S}: \{v_{r,c} : r \in D\} = \{v_{r,c} : r \in \mu_{\delta t}(D)\}$,
rendering all marginal-distribution tests powerless regardless of $|\delta t|$.

% ═════════════════════════════════════════════════════════════════════════════
\section{Semantic Contract Framework}
\label{sec:framework}
% ═════════════════════════════════════════════════════════════════════════════

\subsection{Mutation Taxonomy}

We define five semantic mutation categories, each motivated by a class of
real ETL and pipeline failure modes observed in production ML systems.

\textbf{Scaling mutations (6).}  Scaling errors arise when a unit conversion
is applied incorrectly, a calibration constant is updated in one part of a
pipeline but not another, or a normalisation step is applied twice.
Representative operators include multiplying \texttt{Weekly\_Sales} by 0.9
(10\% systematic undercount), multiplying \texttt{Voltage} by 0.5 (sensor
gain halved after firmware update), and multiplying all price columns by 10
(cents-to-dollars conversion applied twice).

\textbf{Temporal mutations (6).}  Temporal errors destroy the causal and
chronological relationships that time-series ML models depend upon.
Operators include shifting all dates by $\pm 7$ days (weekly cron-job
arithmetic error), lagging stock timestamps by one trading day ($T+1$
settlement date confusion), reversing the entire price series (as-of-date
inversion), and shuffling records within each calendar day (intra-day
ordering destroyed by parallel batch processing).

\textbf{Structural mutations (9).}  Structural errors alter the relational
topology of the dataset---which entities are present, how columns are
related, and whether the row count is consistent with the upstream source.
Operators include swapping Open/Close columns, swapping Voltage/Global
intensity columns, dropping all rows for 5 random stores, dropping a
contiguous 6-hour segment, tiling the entire power dataset 5 times (pipeline
loop bug), and negating \texttt{Global\_reactive\_power} (sign convention
error).

\textbf{Quality mutations (6).}  Quality errors introduce local artefacts
that corrupt individual measurements without changing the overall dataset
structure.  Operators include injecting 100 spikes at $3\sigma$, injecting
50 outliers at $10\sigma$, freezing a 2-hour window at the column mean
(sensor lock-up), zeroing 30\% of volume readings (sparse data after join
bug), duplicating 5\% of rows (ETL double-append), and adding Gaussian
noise at $\sigma = 0.01 \times \sigma_{\mathrm{col}}$.

\textbf{Aggregation mutations (1).}  A single but high-impact operator:
multiplying \texttt{Weekly\_Sales} by 7, representing the most common
aggregation error in retail pipelines---applying a daily summation to data
already aggregated at the weekly level.

Table~\ref{tab:mutations} summarises the complete taxonomy.  In total, 9
mutations target the Walmart dataset, 8 target the stock dataset, and 11
target the power dataset, for 28 mutations in total.

\begin{table*}[!t]
\renewcommand{\arraystretch}{1.15}
\caption{Complete Mutation Taxonomy (28 Operators Across 3 Datasets)}
\label{tab:mutations}
\centering
\small
\begin{tabular}{llll}
\toprule
\textbf{Mutation Name} & \textbf{Dataset} & \textbf{Category} & \textbf{Description} \\
\midrule
\mut{walmart\_\_shift\_dates\_-7}        & Walmart & Temporal     & Shift all dates backward 7 days \\
\mut{walmart\_\_shift\_dates\_7}         & Walmart & Temporal     & Shift all dates forward 7 days \\
\mut{walmart\_\_scale\_sales\_0.9}       & Walmart & Scaling      & Multiply \texttt{Weekly\_Sales} by 0.9 \\
\mut{walmart\_\_scale\_temperature\_2x}  & Walmart & Scaling      & Multiply \texttt{Temperature} by 2 \\
\mut{walmart\_\_toggle\_holiday}         & Walmart & Structural   & Flip \texttt{Holiday\_Flag} $0\leftrightarrow1$ \\
\mut{walmart\_\_missing\_recent\_10pct}  & Walmart & Structural   & Drop most recent 10\% of rows \\
\mut{walmart\_\_missing\_stores\_5}      & Walmart & Structural   & Drop all rows for 5 random stores \\
\mut{walmart\_\_duplicate\_rows\_5pct}   & Walmart & Quality      & Duplicate a random 5\% of rows \\
\mut{walmart\_\_wrong\_aggregation\_x7}  & Walmart & Aggregation  & Multiply \texttt{Weekly\_Sales} by 7 \\
\midrule
\mut{stock\_\_noise\_0.01}               & Stock   & Quality      & Add Gaussian noise $\sigma=0.01\sigma_{\mathrm{col}}$ to prices \\
\mut{stock\_\_zero\_volume\_30pct}       & Stock   & Quality      & Zero out \texttt{Volume} for 30\% of rows \\
\mut{stock\_\_drop\_5pct}                & Stock   & Structural   & Drop 5\% of rows randomly \\
\mut{stock\_\_column\_swap\_open\_close} & Stock   & Structural   & Swap \texttt{Open} and \texttt{Close} columns \\
\mut{stock\_\_shift\_price\_0.5}         & Stock   & Scaling      & Add \$0.50 to all price columns \\
\mut{stock\_\_unit\_conversion\_x10}     & Stock   & Scaling      & Multiply all prices by 10 \\
\mut{stock\_\_timestamp\_lag\_1day}      & Stock   & Temporal     & Lag all timestamps by 1 trading day \\
\mut{stock\_\_reverse\_series}           & Stock   & Temporal     & Reverse the entire price series \\
\midrule
\mut{power\_\_scale\_active\_1.1}              & Power & Scaling    & Multiply \texttt{Global\_active\_power} by 1.1 \\
\mut{power\_\_scale\_voltage\_0.5}             & Power & Scaling    & Multiply \texttt{Voltage} by 0.5 \\
\mut{power\_\_shuffle\_within\_day}            & Power & Temporal   & Shuffle records within each calendar day \\
\mut{power\_\_shift\_by\_1week}                & Power & Temporal   & Shift all timestamps forward 7 days \\
\mut{power\_\_column\_swap\_voltage\_intensity}& Power & Structural & Swap \texttt{Voltage} and \texttt{Global\_intensity} \\
\mut{power\_\_missing\_6hour\_segment}         & Power & Structural & Delete a contiguous 6-hour block \\
\mut{power\_\_data\_repetition\_5x}            & Power & Structural & Tile entire dataset $5\times$ \\
\mut{power\_\_sign\_flip\_reactive}            & Power & Structural & Negate \texttt{Global\_reactive\_power} \\
\mut{power\_\_inject\_spikes\_100\_3.0}        & Power & Quality    & Inject 100 spikes at $3\sigma$ \\
\mut{power\_\_extreme\_outliers\_50}           & Power & Quality    & Inject 50 outliers at $10\sigma$ \\
\mut{power\_\_constant\_segment\_2h}           & Power & Quality    & Freeze 2-hour window at column mean \\
\bottomrule
\end{tabular}
\end{table*}

\subsection{Contract Design Principles}

Each semantic contract implements one or more of the following design
principles, chosen to make detectors robust across datasets and robust to
seed-specific stochastic variation.
Table~\ref{tab:principles} maps each principle to the mutation categories it
targets; the full set ensures no category is detectable by only a single
signal.

\textbf{Alignment-First.}  Before any comparison, the baseline and test
datasets are aligned on their datetime index.  For multi-entity datasets
(e.g., Walmart with 45 stores), alignment is performed by sorting both
frames on the temporal column and truncating to the minimum shared length.
This eliminates spurious differences due to row-count mutations propagating
into ratio computations.

\textbf{Ratio-Based Detection.}  For scaling contracts, we compute the
median of pointwise ratios $b_i / a_i$ between aligned test and baseline
values.  The deviation $|\tilde{r} - 1|$ from the expected ratio of 1.0
is compared against a fractional tolerance.  Using the median rather than
the mean provides robustness to the $\leq$5\% of rows that may differ
legitimately (missing values, outliers).

\textbf{Correlation-Based Structural Detection.}  Column swaps are detected
by computing cross-correlations: $\rho(a^{c_1}, b^{c_2})$ and
$\rho(a^{c_2}, b^{c_1})$, where $a$ and $b$ are aligned baseline and test
arrays.  A cross-correlation advantage $\max(\rho_{12} - \rho_{11},\,
\rho_{21} - \rho_{22}) > \theta$ indicates that test columns align better
with swapped baseline columns than with their own.  A complementary spread
inversion signal---$\rho(\text{spread}_a, \text{spread}_b) < -0.5$---fires
when a column swap inverts the intraday spread, providing a second
independent detection path.

\textbf{ACF-Profile Comparison.}  Temporal mutations are detected by
comparing autocorrelation function profiles at lags $1$ through $24$
(minutes for power data, trading days for stock data).  The mean absolute
difference $\overline{|\mathrm{ACF}_a - \mathrm{ACF}_b|}$ exceeding a
threshold detects both temporal shuffles and timestamp lags.  A
complementary reversal signal compares $\rho(a, b)$ against
$\rho(a, b^{\mathrm{rev}})$: when the series is reversed,
$\rho(a, b^{\mathrm{rev}}) \gg \rho(a, b)$.

\textbf{Spike Isolation.}  Quality contracts for spike injection count the
number of newly appearing spikes in test data relative to baseline:
positions $i$ where $b_i > \bar{a} + k\sigma_a$ but
$a_i \leq \bar{a} + k\sigma_a$.  Using the baseline's own mean and standard
deviation as the threshold makes this detector adaptive to any absolute
signal scale.

\textbf{Entity-Count Validation.}  For multi-entity datasets, dropping all
rows for a subset of entities is detected by comparing the set of unique
store identifiers.  A single missing entity (threshold: $\geq 1$ missing)
triggers detection with confidence proportional to the fraction of the
entity vocabulary missing.

\textbf{Row-Count Ratio.}  Duplication and repetition mutations are detected
by the fractional excess in test row count:
$\texttt{extra\_frac} = (n_{\mathrm{test}} - n_{\mathrm{base}}) /
n_{\mathrm{base}}$.  This is the most direct and reliable detector for
duplicate rows, since duplicated rows are statistically indistinguishable
from originals.

\textbf{Date-Range Comparison.}  Timestamp mutations are detected by
comparing the median date of test versus baseline.  If all records are
shifted by $\delta$ days, the median shifts by exactly $\delta$ days.
This is robust to the multi-store structure of the Walmart dataset, where
no per-row temporal comparison is possible.

\begin{table}[!t]
\renewcommand{\arraystretch}{1.15}
\caption{Contract Design Principles and Target Categories}
\label{tab:principles}
\centering
\small
\begin{tabular}{lp{3.8cm}}
\toprule
\textbf{Principle} & \textbf{Targets} \\
\midrule
Alignment-First            & All categories \\
Ratio-Based Detection      & Scaling, Aggregation \\
Correlation-Based Struct.\ & Structural (column swaps) \\
ACF-Profile Comparison     & Temporal \\
Spike Isolation            & Quality (spikes, outliers) \\
Entity-Count Validation    & Structural (missing entities) \\
Row-Count Ratio            & Quality (duplicates), Structural \\
Date-Range Comparison      & Temporal (shifts, lags) \\
\bottomrule
\end{tabular}
\end{table}

\subsection{Dataset-Specific Contract Deployment}

For the Walmart dataset, we deploy 7 contracts: shift-dates, scale-sales,
toggle-holiday, scale-temperature, missing-recent, missing-stores, and
duplicate-rows.  For the stock dataset, we deploy 8 contracts: noise,
drop-fraction, shift-price, column-swap (Open/Close pair), unit-conversion,
temporal-ACF (with reversal signal), zero-volume, and timestamp-lag.  For
the power dataset, we deploy 10 contracts: scale-active-power,
scale-voltage, inject-spikes, shuffle-within-day, column-swap
(Voltage/Global\_intensity pair), missing-segment, data-repetition,
sign-flip-reactive, temporal-ACF, and constant-segment.

Each contract returns a binary \texttt{detected} flag and a normalised
\texttt{confidence} score in $[0,1]$.  The semantic validator aggregates
over all contracts for the relevant dataset: a mutation is flagged as
detected if at least one contract fires with confidence $\geq 0.10$.
This low threshold maximises recall; the FPR = 0.000 result confirms that
the contracts do not fire spuriously on clean data.

\subsection{Semantic Precision}

Each detector is assigned to exactly one of the five taxonomy categories
(e.g., \texttt{scale\_sales} $\rightarrow$ \textsc{scaling},
\texttt{temporal\_acf} $\rightarrow$ \textsc{temporal}).  Each mutation has
a known ground-truth category.  Semantic precision (Eq.~\ref{eq:sp})
measures whether the correct category of detector fires, not merely whether
any detector fires.  A detection where a \textsc{scaling} contract fires on
a \textsc{structural} mutation is counted as a semantic precision failure,
even though the mutation is correctly flagged as detected.  This metric
provides a stronger claim: the framework does not merely detect anomalies,
it correctly diagnoses their category.
One implementation note: the sole Aggregation mutation
(\mut{walmart\_\_wrong\_aggregation\_x7}) is detected by the
\texttt{scale\_sales} contract, which is categorised as Scaling.  Because
a $7\times$ sales inflation is mechanistically identical to a scaling error
from the contract's perspective, the semantic precision computation maps
Aggregation mutations to the Scaling detector category for the purposes of
precision scoring.  This is disclosed explicitly so that the 100\%
Aggregation precision figure is not overstated: it reflects the fact that
the correct ratio-based signal fires, not that a dedicated aggregation
contract exists.


\section{Experimental Setup}
\label{sec:setup}


\subsection{Datasets}

We use three publicly available tabular datasets representing distinct
domains and data modalities.

\textbf{Walmart Retail Sales \cite{walmart2014kaggle}.}  The Walmart Store
Sales Forecasting dataset contains 6,435 rows across 45 store-week
combinations.  Key columns include \texttt{Weekly\_Sales},
\texttt{Temperature}, \texttt{Holiday\_Flag}, \texttt{Fuel\_Price},
\texttt{CPI}, \texttt{Unemployment}, \texttt{Store}, and \texttt{Date}.
This dataset is representative of enterprise retail ML pipelines, where
aggregation errors and holiday-flag inversions are common ETL failure modes.

\textbf{AAPL Stock Prices (Yahoo Finance, 5-year daily OHLCV).}  1,259
rows of daily Open, High, Low, Close, Adjusted Close, and Volume for Apple
Inc.\ stock over the period 2019--2024.  This dataset stresses temporal
ordering constraints and inter-column semantic relationships (e.g.,
High $\geq$ Close $\geq$ Low on any given day).

\textbf{UCI Household Power Consumption \cite{uci2013}.}  A 996,010-row
sample of per-minute household electrical measurements over four years, with
columns \texttt{Global\_active\_power}, \texttt{Global\_reactive\_power},
\texttt{Voltage}, \texttt{Global\_intensity}, and three sub-metering
channels.  This dataset stresses quality contracts (spike detection, sensor
freeze) and temporal structure (intra-day autocorrelation patterns).

\subsection{Mutation Generation and Reproducibility}

Each of the 28 mutation operators is parameterised by a random seed.  We
run all experiments with seeds $\{42, 123, 456, 789\}$, producing
$28 \times 4 = 112$ mutated trials.  Seeds are used exclusively for
stochastic operators (e.g., selecting which 5\% of rows to duplicate, which
stores to drop, where to place spikes); purely deterministic operators
(e.g., $\times 0.9$ scaling, $+7$-day date shift) are unaffected by seed
choice.  All code, datasets, seeds, and output CSVs will be released in a
public repository upon acceptance to ensure full reproducibility.  The
complete pipeline can be executed with a single command and reproduces all
reported figures exactly.

\subsection{Baseline Validators}

\textbf{KS Drift (original baseline).}  A per-column two-sample
Kolmogorov-Smirnov test \cite{kolmogorov1933} with Bonferroni correction.
A mutation is detected if any column yields $p < 0.01 / n_{\mathrm{cols}}$.
This is the standard production drift detector.

\textbf{Ensemble Statistical.}  A five-test ensemble combining KS test,
Jensen-Shannon divergence \cite{lin1991jsd} on binned empirical
distributions, Population Stability Index \cite{yurdakul2018psi}, Frobenius
norm of Pearson inter-column correlation matrix difference, and mean ACF
difference over lags 1--20.  Detection uses any-one (OR) voting: a mutation
is detected if at least one test fires.  This baseline is designed to be the
strongest possible purely statistical detector.

\textbf{Great Expectations Comprehensive \cite{greatexpectations2023}.}  A
calibrated expectation suite mimicking production GE deployments.
Eight expectation types are applied: row count within $\pm$15\% of
baseline, column mean within $\pm$20\% (or $\pm$30\% for stock prices),
column standard deviation within [50\%, 200\%] of baseline std, null rate
increase $\leq$10\,pp, 99\% of values within an expanded baseline
$[q_{0.01}, q_{0.99}]$ range, zero-fraction increase $\leq$20\,pp,
categorical completeness ($\leq$10\% unique values missing), and sign
distribution for reactive power.  Tolerances are deliberately generous to
match production usage and minimise false alerts.

\textbf{Semantic Validator (proposed).}  The framework described in
Section~\ref{sec:framework}, with confidence threshold 0.10.

\subsection{Evaluation Protocol and Statistical Methods}

All four validators are applied to each of the 112 mutated trials and to
each of the 12 clean baseline trials (3 datasets $\times$ 4 seeds).
We report:

\begin{itemize}
\item \textbf{Pooled Detection Rate} with 95\% Wilson score confidence
      intervals \cite{wilson1927}.  The Wilson interval is preferred over
      the Wald interval for proportions that may be near 0 or 1, as it
      maintains nominal coverage for small to moderate $n$.

\item \textbf{McNemar exact test} \cite{mcnemar1947} for pairwise validator
      comparisons.  The exact binomial version (rather than asymptotic
      $\chi^2$) is appropriate here because the number of discordant pairs
      is small relative to the sample size.  Significance threshold:
      $\alpha = 0.05$.

\item \textbf{Two-proportion z-test} for the development-vs-held-out
      generalisation gap.

\item \textbf{Semantic precision} (Eq.~\ref{eq:sp}).
\end{itemize}

Wilson score intervals and McNemar tests were implemented 
using the \texttt{statsmodels} statistical library~\cite{seabold2010statsmodels}.

We report per-seed detection rates alongside pooled rates to verify that
confidence intervals are not artificially tightened by treating correlated
per-seed observations as independent.  The standard deviation across seeds
is 0.0\% for the semantic validator, confirming that the pooled CI reflects
genuine consistency rather than artefact.

% ═════════════════════════════════════════════════════════════════════════════
\section{Results}
\label{sec:results}
% ═════════════════════════════════════════════════════════════════════════════

\subsection{Main Detection Results}

Table~\ref{tab:main} presents the primary comparison.  The semantic
validator detects 108 of 112 mutated trials (96.4\%), more than doubling
the detection rate of the best statistical baseline (Ensemble: 49.1\%).
The $7.3\times$ advantage over the KS-only baseline (42.0\%) and the
$2.2\times$ advantage over the ensemble baseline are consistent across all
four seeds.  All four validators achieve FPR = 0.000, establishing that the
improvements in detection rate are not obtained by trading off precision.

Notably, 12 of the 28 mutations are detected at 100\% by the semantic
validator while achieving 0\% detection across \emph{all three} baselines
simultaneously---including \mut{power\_\_shift\_by\_1week},
\mut{stock\_\_column\_swap\_open\_close}, \mut{stock\_\_shift\_price\_0.5},
\mut{stock\_\_timestamp\_lag\_1day}, \mut{walmart\_\_duplicate\_rows\_5pct},
\mut{walmart\_\_shift\_dates\_-7}, and \mut{walmart\_\_shift\_dates\_7},
among others.  These represent the mutation types for which statistical
monitoring provides no protection whatsoever.

\begin{figure}[!t]
\centering
\includegraphics[width=\columnwidth]{fig1_detection_rates.pdf}
\caption{Pooled detection rate (95\% Wilson CI) for all four 
validators across 112 mutated trials. Error bars show Wilson 
score confidence intervals. The semantic validator achieves 
96.4\%, more than doubling the best statistical baseline.}
\label{fig:dr_comparison}
\end{figure}

\begin{table}[!t]
\renewcommand{\arraystretch}{1.25}
\caption{Detection Rate Comparison ($28 \times 4 = 112$ Trials)}
\label{tab:main}
\centering
\small
\begin{tabular}{lccc}
\toprule
\textbf{Validator} & \textbf{Det.} & \textbf{DR (95\% Wilson CI)} & \textbf{FPR} \\
\midrule
Semantic Validator (ours) & 108/112 & \textbf{96.4\% [91.2\%, 98.6\%]} & 0.000 \\
Ensemble Statistical      & 55/112  & 49.1\% [40.0\%, 58.2\%]          & 0.000 \\
KS Drift (original)       & 47/112  & 42.0\% [33.2\%, 51.2\%]          & 0.000 \\
GE Comprehensive          & 48/112  & 42.9\% [34.1\%, 52.1\%]          & 0.000 \\
\bottomrule
\multicolumn{4}{l}{\footnotesize FPR denominator: 12 clean trials (3 datasets $\times$ 4 seeds).}
\end{tabular}
\end{table}

\subsection{Statistical Significance}

Table~\ref{tab:mcnemar} presents McNemar exact test results for all
pairwise comparisons between the semantic validator and each baseline.  The
number of discordant pairs (mutations detected by Semantic but not Baseline,
or vice versa) is large in all cases, yielding $p < 0.0001$ for every
comparison.  The advantage of the semantic validator over all three baselines
is statistically overwhelming and cannot be attributed to chance.

\begin{table}[!t]
\renewcommand{\arraystretch}{1.25}
\caption{McNemar Exact Test Results ($\alpha = 0.05$)}
\label{tab:mcnemar}
\centering
\begin{tabular}{lcc}
\toprule
\textbf{Comparison} & \textbf{$p$-value} & \textbf{Significant} \\
\midrule
Semantic vs.\ Ensemble Statistical & $< 0.0001$ & Yes \\
Semantic vs.\ KS Drift (original)  & $< 0.0001$ & Yes \\
Semantic vs.\ GE Comprehensive     & $< 0.0001$ & Yes \\
\bottomrule
\end{tabular}
\end{table}

\subsection{Per-Category Breakdown}

Table~\ref{tab:category} disaggregates detection rates by mutation category
for the semantic validator.  Four of the five categories achieve 100\%
detection.  The single below-perfect category is Quality at 83.3\%
(20/24), attributable entirely to a single mutation
(\mut{stock\_\_noise\_0.01}) across all four seeds.  This mutation adds
Gaussian noise at $\sigma = 0.01 \times \sigma_{\mathrm{col}}$---a
perturbation whose magnitude is sub-1\% of natural price variability---and
is undetectable by any of the four validators.  We discuss this as a
documented detection limit in Section~\ref{sec:discussion}.

\begin{figure}[!t]
\centering
\includegraphics[width=\columnwidth]{fig2_category_heatmap.pdf}
\caption{Per-category detection rate (\%) for all four validators. 
The semantic validator achieves 100\% in four of five categories. 
The KS baseline detects 0\% of Temporal mutations; Ensemble and GE 
detect only 33\%, compared with 100\% for the semantic validator.}
\label{fig:category_heatmap}
\end{figure}

\begin{table}[!t]
\renewcommand{\arraystretch}{1.25}
\caption{Per-Category Detection Rate (Semantic Validator)}
\label{tab:category}
\centering
\begin{tabular}{lccl}
\toprule
\textbf{Category} & \textbf{Det.} & \textbf{Total} &
  \textbf{DR (95\% Wilson CI)} \\
\midrule
Aggregation & 4   & 4   & 100.0\% [51.0\%, 100.0\%] \\
Quality     & 20  & 24  & \phantom{0}83.3\% [64.1\%, \phantom{0}93.3\%] \\
Scaling     & 24  & 24  & 100.0\% [86.2\%, 100.0\%] \\
Structural  & 36  & 36  & 100.0\% [90.4\%, 100.0\%] \\
Temporal    & 24  & 24  & 100.0\% [86.2\%, 100.0\%] \\
\midrule
\textbf{Overall} & \textbf{108} & \textbf{112} &
\textbf{96.4\% [91.2\%, 98.6\%]} \\
\bottomrule
\end{tabular}
\end{table}

\subsection{Semantic Precision}

Table~\ref{tab:precision} reports per-category semantic precision.  The overall semantic precision of 96.3\% (104/108) demonstrates that the
framework not only detects mutations but correctly diagnoses their category
in the vast majority of cases.  Four categories achieve 100\% precision.
The sole precision shortfall occurs in the Structural category (88.9\%,
32/36), and is entirely attributable to one mutation---\mut{stock\_\_drop\_5pct}---across
all four seeds.  Dropping 5\% of rows shifts the apparent price level
and truncates the timestamp range by a small amount, causing the
\texttt{shift\_price} (Scaling) and \texttt{timestamp\_lag} (Temporal)
contracts to fire before the \texttt{drop\_fraction} (Structural) contract
reaches its threshold.  The mutation is correctly \emph{detected} in all
four seed runs; the precision failure reflects incidental cross-category
detector firing, not a missed mutation.  No Quality precision failures
occur: all 20 detected Quality mutations are attributed to the correct
category.

\begin{table}[!t]
\renewcommand{\arraystretch}{1.25}
\caption{Semantic Precision by Category}
\label{tab:precision}
\centering
\begin{tabular}{lccr}
\toprule
\textbf{Category} & \textbf{Detected} & \textbf{Sem.\ Matches} &
  \textbf{Precision} \\
\midrule
Aggregation & 4   & 4   & 100.0\% \\
Quality     & 20  & 20  & 100.0\% \\
Scaling     & 24  & 24  & 100.0\% \\
Structural  & 36  & 32  & \phantom{0}88.9\% \\
Temporal    & 24  & 24  & 100.0\% \\
\midrule
\textbf{Overall} & \textbf{108} & \textbf{104} & \textbf{96.3\%} \\
\bottomrule
\end{tabular}
\end{table}

\subsection{Seed Stability}

Table~\ref{tab:seeds} reports per-seed detection rates for all four
validators.  The semantic validator is perfectly stable at 96.4\% across
all seeds, yielding standard deviation 0.0\%.  This eliminates the
possibility that the pooled confidence interval is artificially tightened
by within-mutation correlation: if the four observations per mutation were
strongly correlated, we would expect seed-to-seed variance; the absence of
any variance confirms the pooled CI is conservative, not liberal.  The ensemble baseline shows minor seed-to-seed variation 
(range: 46.4\%--50.0\%) reflecting the continuous nature 
of its component test scores.

\begin{table}[!t]
\renewcommand{\arraystretch}{1.15}
\caption{Per-Seed Detection Rates (\%)}
\label{tab:seeds}
\centering
\begin{tabular}{lccccc}
\toprule
\textbf{Validator} & \textbf{S=42} & \textbf{S=123} &
  \textbf{S=456} & \textbf{S=789} & \textbf{Std} \\
\midrule
Semantic Validator & 96.4 & 96.4 & 96.4 & 96.4 & 0.0 \\
Ensemble Stat.     & 46.4 & 50.0 & 50.0 & 50.0 & --- \\
KS Drift           & 39.3 & 42.9 & 42.9 & 42.9 & --- \\
GE Comprehensive   & 42.9 & 42.9 & 42.9 & 42.9 & 0.0 \\
\bottomrule
\end{tabular}
\end{table}

\subsection{Held-Out Evaluation}

To assess whether the contracts generalise beyond the specific mutation
types observed during iterative development, we partition the 28 mutations
into a \emph{development set} of 19 mutations (whose failure modes were
observed during detector tuning) and a \emph{held-out set} of 9 mutations
(for which no detector was written or revised after observing a failure).

Table~\ref{tab:heldout} presents results.  On the development set, the
semantic validator achieves 100.0\% (95\% CI: 95.2\%--100.0\%) with all
five categories perfect.  On the held-out set, it achieves 88.9\% (95\% CI:
74.7\%--95.6\%), a generalisation gap of $+11.1$ percentage points
(two-proportion z-test: $z = 2.959$, $p = 0.003$).  The held-out detection
rate exceeds all three baselines by $>$36 percentage points.

% FIX 7: Table VIII previously had multi-line text inside narrow cells
%         causing CI values to wrap mid-number in two-column layout.
%         Reformatted to use a p-column with explicit width.
\begin{table}[!t]
\renewcommand{\arraystretch}{1.25}
\caption{Held-Out Evaluation Results}
\label{tab:heldout}
\centering
\small
\begin{tabular}{lcccc}
\toprule
\textbf{Set} & \textbf{Sem DR (95\% CI)} & \textbf{Ens} & \textbf{KS} & \textbf{GE} \\ 
\midrule
Dev.\ (19)   & \scriptsize 100.0\% {[95.2, 100.0]} & 47.4\% & 42.1\% & 42.1\% \\
Held-Out (9) & \scriptsize 88.9\% {[74.7,\;95.6\phantom{0}]} & 52.8\% & 41.7\% & 44.4\% \\
All (pooled) & \scriptsize 96.4\% {[91.2,\;98.6\phantom{0}]} & 49.1\% & 42.0\% & 42.9\% \\
\midrule
\multicolumn{5}{l}{\footnotesize 95\% Wilson CI bounds shown in brackets.}\\
\multicolumn{5}{l}{\footnotesize Gap: $+11.1$\,pp ($z=2.959$, $p=0.003$).}\\
\bottomrule
\end{tabular}
\end{table}

Per-mutation held-out results reveal the structure of the gap.  Eight of the
nine held-out mutations are detected at 100\% by the semantic validator; the
single miss across all four seeds is \mut{stock\_\_noise\_0.01} (0\% for
all validators).  The remaining eight include the particularly challenging
\mut{stock\_\_timestamp\_lag\_1day} (detected at 100\% by semantic contracts
via date-range comparison, at 0\% by Ensemble and KS) and
\mut{stock\_\_reverse\_series} (100\% by semantic and Ensemble via ACF
reversal signal, 0\% by KS and GE).  One caveat applies to
\mut{stock\_\_drop\_5pct}: it is correctly detected at 100\% but via
incidental Scaling and Temporal detector signals rather than its targeted
Structural contract, as discussed in Section~\ref{sec:results} (Semantic
Precision).  This does not affect the detection rate but is reflected in
the Structural category precision of 88.9\%.

\subsection{Missed Mutation Analysis and False Positive Analysis}

The single mutation that consistently evades detection across all four
validators is \mut{stock\_\_noise\_0.01}, which injects Gaussian noise at
$\sigma_{\mathrm{noise}} = 0.01 \times \sigma_{\mathrm{col}}$ to all price
columns.  This perturbation is sub-1\% of the natural daily price
variability and is therefore genuinely indistinguishable from benign
measurement noise.  Lowering the noise detection threshold to catch this
mutation would generate false positives on clean market data.  This case
represents a \emph{fundamental detection limit} of threshold-based
validation against imperceptibly small perturbations, not a weakness of the
semantic contract paradigm specifically.

% FIX 8: Original had "n>10,000n > 10,000 n>10,000" — triplicated
%         text/math artefact.  Replaced with single clean math expression.
A technical note: scipy's \texttt{ks\_2samp} function emits 
a runtime warning on the 996{,}010-row power dataset, switching 
from exact to asymptotic $p$-value computation. This fallback 
is mathematically valid; the asymptotic approximation is known 
to be accurate for $n > 10{,}000$~\cite{kolmogorov1933} and 
does not affect any reported result.

Excluding \mut{stock\_\_noise\_0.01}, the held-out detection rate is
100.0\%, and the generalisation gap collapses to 0 percentage points.

The false positive rate is 0.000 across all 12 clean baseline evaluations
and all four validators.  In a production context, this means the framework
generates zero unnecessary alerts on clean data---a prerequisite for
adoption in operational environments where false alarms erode trust and
lead to validator bypass.

% ═════════════════════════════════════════════════════════════════════════════
\section{Discussion}
\label{sec:discussion}
% ═════════════════════════════════════════════════════════════════════════════

\subsection{Statistical versus Semantic Validation: The Core Argument}

The results establish, both theoretically (Section~\ref{sec:formulation})
and empirically (Section~\ref{sec:results}), that statistical methods are
structurally insufficient for semantic mutation detection.  The insight is
not that statistical tests are poorly implemented or inadequately
parameterised---our ensemble baseline combines five state-of-the-art tests
with any-one voting, representing the strongest conceivable purely
statistical detector---but that statistical tests fundamentally operate on
the wrong representation of the data.

A statistical test asks: \emph{do these two samples look like they came
from the same distribution?}  A semantic contract asks: \emph{do these two
datasets maintain the relational invariants that the domain requires?}
These are different questions.  Column swaps, timestamp shifts, aggregation
errors, and sign flips all answer ``yes'' to the first question and ``no''
to the second.

An instructive comparison is the GE baseline, which achieves only 42.9\%
despite incorporating eight categories of expectations including categorical
completeness and value range checks.  Inspection of per-mutation results
shows that GE and KS share \emph{identical blind spots} on temporal
mutations and subtle structural anomalies.  Despite being a richer
framework, GE fails for the same structural reason as KS: its expectations
are defined over the marginal properties of columns, not over the
semantic relationships between columns and the domain they represent.
Adding more statistical expectations to a GE suite does not close the
gap; only expectations that encode domain meaning do.

% FIX 9: Original had "powershif tby1week" — garbled mutation name caused
%         by \texttt hyphenation in two-column layout.
%         Replaced with \mut{} macro (= \mbox{\texttt{...}}) throughout.
Concretely, the KS baseline achieves \textbf{0\% detection on all six
Temporal mutations} across all seeds, and GE detects only two of six
(33\%).  The three held-out temporal mutations---\mut{power\_\_shift\_by\_1week},
\mut{stock\_\_timestamp\_lag\_1day}, and \mut{stock\_\_reverse\_series}---are
detected at 0\% by both GE and KS, while the semantic validator detects
all six Temporal mutations at 100\%.  This is not a
matter of threshold calibration---no threshold adjustment to a
marginal-distribution test can detect a timestamp shift that leaves all
marginal distributions identical.

\subsection{Limitations}

We report the following limitations honestly, as they determine the scope
of valid inferences from this work.

\textbf{L1\,---\,Domain specificity.}  The contracts presented here are
hand-crafted for three specific datasets.  Writing effective contracts for a
new dataset requires domain expertise to identify the invariants that the
domain mandates.  The framework is not plug-and-play; it requires upfront
investment from a domain expert.  We view this as a feature rather than a
bug---domain knowledge is precisely what statistical tools lack---but it
does impose a deployment cost.

\textbf{L2\,---\,Approximate held-out partition.}  The development/held-out
split is based on development history: which mutation types were observed
during iterative contract refinement.  This history cannot be fully audited.
Broad design choices (e.g., the decision to include ACF-based contracts) may
have been influenced by awareness of all mutation types.  The held-out
results should therefore be interpreted as a partial, not strict,
generalisation test.

\textbf{L3\,---\,Statistically significant generalisation gap.}  The
$+11.1$ percentage-point gap between development (100.0\%) and held-out
(88.9\%) detection rates is statistically significant ($p = 0.003$).
Although this gap is entirely attributable to a single documented hard-limit
case (\mut{stock\_\_noise\_0.01}), this attribution should not be used to
dismiss the gap.  Future work should investigate held-out evaluation on a
fully disjoint dataset with no development-phase exposure.

\textbf{L4\,---\,Confidence interval dependence.}  The pooled detection
rate treats the $4 \times 28 = 112$ trials as independent observations.  In
reality, the four per-mutation observations are correlated (same mutation,
different seed).  The per-seed standard deviation of 0.0\% is the strongest
possible evidence that this correlation does not bias the pooled CI, but the
methodological caveat should be noted.

\textbf{L5\,---\,Threshold sensitivity.}  Contract thresholds (e.g., the
3\% tolerance for the scale-sales detector, the $-0.70$ Pearson threshold
for the sign-flip detector) were set by domain knowledge and iterative
refinement on the development set.  No systematic sensitivity analysis was
performed.  Results may degrade for datasets with unusual signal
characteristics that place clean data close to these thresholds.

\textbf{L6\,---\,Scope of datasets.}  Three datasets spanning three domains
(retail, finance, energy) provide diversity, but do not constitute an
exhaustive evaluation.  The framework has not been tested on NLP features
stored as tabular embeddings, image metadata, or non-temporal tabular data.

\subsection{Threats to Validity}

\textbf{Internal validity.}  Contract thresholds are the primary internal
validity threat.  A threshold that was tuned with (even indirect) knowledge
of the held-out mutations could inflate generalisation results.  We have
mitigated this by making the partition explicit and by reporting the gap
honestly.  A controlled threshold-sensitivity experiment is a recommended
follow-up.

\textbf{External validity.}  Results may not generalise to domains with
higher natural variability (e.g., social media data, sensor networks with
legitimate high-frequency spikes).  In such domains, the contracts would
require recalibration, which may raise FPR.  The three datasets were chosen
to represent distinct variability regimes (low-variance weekly retail;
medium-variance daily equity; high-variance minute-level energy), providing
moderate confidence in the generality of the design principles if not
the specific thresholds.

% ═════════════════════════════════════════════════════════════════════════════
\section{Conclusion}
\label{sec:conclusion}
% ═════════════════════════════════════════════════════════════════════════════

This paper introduces semantic contract-based validation as a principled
solution to a class of data pipeline failure---the semantic mutation---that
is structurally undetectable by existing statistical monitoring tools.  We
formalise the semantic mutation, construct a five-category taxonomy of 28
real-world ETL failure patterns, implement 25 domain-aware contracts, and
demonstrate empirically that the resulting validator achieves 96.4\%
detection rate at 0.000 false-positive rate---more than doubling the
performance of the best available statistical baseline---with statistical
significance at $p < 0.0001$ in all pairwise McNemar tests.  Semantic
precision of 96.3\% confirms that the framework diagnoses mutation types
correctly, not merely flags anomalies.

The practical implication for ML data teams is direct: schema validators
and drift detectors are necessary but not sufficient for data quality
assurance.  A third layer---semantic contracts---is required to catch the
mutations that preserve syntactic and distributional properties while
violating the domain meaning of the data.  Implementing this layer requires
domain expertise but, once written, contracts can be executed in
milliseconds per incoming batch and integrated into any CI/CD pipeline gate.

Three directions warrant future investigation.  First, the framework should
be extended to additional datasets and modalities, including natural language
feature stores and time-series with irregular sampling, to characterise the
conditions under which hand-crafted contracts can be replaced by learned
invariants.  Second, automatic contract generation from data dictionaries or
schema annotations---analogous to constraint suggestion in
Deequ \cite{schelter2018deequ}---would reduce the domain-expert burden and
enable broader adoption.  Third, integration with MLOps platforms (MLflow,
Kubeflow, Vertex AI) as a pre-training validation gate would enable
continuous semantic validation in production, transforming the framework
from a research prototype into a production-grade data reliability
engineering tool.

% ── References ───────────────────────────────────────────────────────────────
\begin{thebibliography}{99}

\bibitem{sculley2015}
D.~Sculley, G.~Holt, D.~Golovin, E.~Davydov, T.~Phillips, D.~Ebner,
V.~Chaudhary, M.~Young, J.-F.~Crespo, and D.~Dennison,
``Hidden technical debt in machine learning systems,''
in \emph{Advances in Neural Information Processing Systems (NIPS)},
vol.~28, 2015, pp.~2503--2511.

\bibitem{polyzotis2018}
N.~Polyzotis, S.~Roy, S.~E.~Whang, and M.~Zinkevich,
``Data lifecycle challenges in production machine learning: A survey,''
\emph{ACM SIGMOD Record}, vol.~47, no.~2, pp.~17--28, 2018.

\bibitem{greatexpectations2023}
\emph{Great Expectations: Always Know What to Expect From Your Data},
Superconductive Inc., 2023. [Online].
Available: \url{https://greatexpectations.io}

\bibitem{kolmogorov1933}
A.~N.~Kolmogorov,
``Sulla determinazione empirica di una legge di distribuzione,''
\emph{Giornale dell'Istituto Italiano degli Attuari}, vol.~4,
pp.~83--91, 1933.

\bibitem{breck2017}
E.~Breck, S.~Cai, E.~Nielsen, M.~Salib, and D.~Sculley,
``The ML test score: A rubric for ML production readiness and
technical debt reduction,''
in \emph{Proc.\ IEEE Int.\ Conf.\ Big Data}, 2017, pp.~1123--1132.

\bibitem{schelter2018automating}
S.~Schelter, D.~Lange, P.~Schmidt, M.~Celikel, F.~Biessmann,
and A.~Grafberger,
``Automating large-scale data quality verification,''
\emph{Proc.\ VLDB Endowment}, vol.~11, no.~12, pp.~1781--1794, 2018.

\bibitem{schelter2018deequ}
S.~Schelter, J.-H.~Lange, P.~Schmidt, M.~Celikel, F.~Biessmann,
and A.~Grafberger,
``Deequ: Data quality validation for machine learning pipelines,''
in \emph{NeurIPS Workshop on Systems for Machine Learning}, 2018.

\bibitem{baylor2017tfx}
D.~Baylor, E.~Breck, H.-T.~Cheng, N.~Fiedel, C.~Y.~Foo, Z.~Haque,
S.~Haykal, M.~Ispir, V.~Jain, L.~Koc, C.~Y.~Koo, L.~Lew, C.~Mewald,
A.~N.~Modi, N.~Polyzotis, S.~Ramesh, S.~Roy, S.~E.~Whang, M.~Wicke,
J.~Wilkiewicz, X.~Zhang, and M.~Zinkevich,
``TFX: A TensorFlow-based production-scale machine learning platform,''
in \emph{Proc.\ ACM SIGKDD Int.\ Conf.\ Knowledge Discovery and
Data Mining}, 2017, pp.~1387--1395.

\bibitem{demillo1978}
R.~A.~DeMillo, R.~J.~Lipton, and F.~G.~Sayward,
``Hints on test data selection: Help for the practicing programmer,''
\emph{Computer}, vol.~11, no.~4, pp.~34--41, Apr.\ 1978.

\bibitem{jia2010}
Y.~Jia and M.~Harman,
``An analysis and survey of the development of mutation testing,''
\emph{IEEE Trans.\ Software Engineering}, vol.~37, no.~5,
pp.~649--678, Sep.\ 2010.

\bibitem{ma2018deepmutation}
L.~Ma, F.~Zhang, J.~Sun, M.~Xue, B.~Li, F.~Juefei-Xu, C.~Xie,
L.~Li, Y.~Liu, J.~Zhao, and Y.~Wang,
``DeepMutation: Mutation testing of deep learning systems,''
in \emph{Proc.\ IEEE 29th Int.\ Symp.\ Software Reliability
Engineering (ISSRE)}, 2018, pp.~100--111.

\bibitem{humbatova2021deepcrime}
N.~Humbatova, G.~Jahangirova, and P.~Tonella,
``DeepCrime: Mutation testing of deep learning systems based on
real faults,''
in \emph{Proc.\ 30th ACM SIGSOFT Int.\ Symp.\ Software Testing
and Analysis (ISSTA)}, 2021, pp.~67--78.

\bibitem{lin1991jsd}
J.~Lin,
``Divergence measures based on the Shannon entropy,''
\emph{IEEE Trans.\ Information Theory}, vol.~37, no.~1,
pp.~145--151, Jan.\ 1991.

\bibitem{yurdakul2018psi}
M.~Yurdakul,
``Statistical properties of Population Stability Index,''
\emph{Western Michigan University, Dept.\ of Statistics}, 2018.

\bibitem{wilson1927}
E.~B.~Wilson,
``Probable inference, the law of succession, and statistical inference,''
\emph{J.\ American Statistical Association}, vol.~22, no.~158,
pp.~209--212, 1927.

\bibitem{mcnemar1947}
Q.~McNemar,
``Note on the sampling error of the difference between correlated
proportions or percentages,''
\emph{Psychometrika}, vol.~12, no.~2, pp.~153--157, 1947.

\bibitem{uci2013}
M.~Lichman,
``UCI Machine Learning Repository,''
University of California, Irvine, School of Information and
Computer Sciences, 2013. [Online].
Available: \url{http://archive.ics.uci.edu/ml}

\bibitem{walmart2014kaggle}
\emph{Walmart Store Sales Forecasting Dataset},
Kaggle, 2014. [Online].
Available: \url{https://www.kaggle.com/c/walmart-recruiting-store-sales-forecasting}

\bibitem{shankar2022operationalizing}
S.~Shankar, R.~Roelofs, H.~Sagara, A.~Namgung, A.~Kumar, and B.~Recht,
``Operationalizing machine learning: An interview study,''
\emph{arXiv preprint arXiv:2209.09125}, 2022.

\bibitem{reis2022}
J.~Reis and M.~Housley,
\emph{Fundamentals of Data Engineering: Plan and Build Robust Data Systems}.
O'Reilly Media, 2022.

\bibitem{seabold2010statsmodels}
S.~Seabold and J.~Perktold,
``Statsmodels: Econometric and statistical modeling with Python,''
in \emph{Proc.\ 9th Python in Science Conf.\ (SciPy)}, 2010,
pp.~92--96.

\bibitem{makinen2021mlops}
S.~M\"{a}kinen, H.~Skogstr\"{o}m, E.~Laaksonen, and T.~Mikkonen,
``Who needs MLOps: What data scientists seek to accomplish and how
can MLOps help?,''
in \emph{Proc.\ IEEE/ACM 1st Workshop on AI Engineering---Software
Engineering for AI (WAIN)}, 2021, pp.~109--112.

\end{thebibliography}

\end{document}